\section{Итоги работы}

\begin{frame}
    \frametitle{Научная новизна}
    \small
    \begin{enumerate}
        \item Предложена модель распределения вычислительных задач AMR между периферийными узлами с учётом дедлайнов, приоритетов, задержек, очередей и внешнего эффекта назначения.
        \item Обосновано использование проектирования экономических механизмов как аппарата внутренней координации общего вычислительного ресурса.
        \item Разработан VCG-подобный аукционный слой с интерпретацией платежей как внутренних теневых цен.
        \item Формализовано адаптивное управление режимами edge-узлов как Dec-POMDP.
        \item Разработаны AMPLE-AMR и C-AMPLE-AMR, объединяющие аукционное распределение, QMIX и кластеризацию инфраструктуры.
    \end{enumerate}
\end{frame}

\begin{frame}
    \frametitle{Практическая значимость}
    \small
    \begin{itemize}
        \item Подход применим при проектировании распределения вычислительных задач в роботизированных складах, производственной и терминальной логистике, больничной логистике и других локальных AMR-средах.
        \item Программный прототип позволяет воспроизводимо сравнивать методы распределения по благосостоянию, задержке, отказам, загрузке узлов и накладным расходам управления.
        \item Экспериментальный контур автоматически формирует CSV, графики и \LaTeX-таблицы для диссертации.
        \item Получены практические правила: адаптивные режимы полезны уже на средних масштабах; кластерное распределение следует включать только при существенном росте накладных расходов глобального раунда.
    \end{itemize}
\end{frame}

\begin{frame}
    \frametitle{Основные результаты к аттестации}
    \small
    \begin{enumerate}
        \item Завершён текст диссертации и согласована структура исследования.
        \item Разработаны математическая модель, аукционный слой, MARL-формализация и гибридный метод AMPLE-AMR.
        \item Разработана масштабируемая версия C-AMPLE-AMR на основе кластеризации графа инфраструктуры.
        \item Реализован программный прототип и набор сценариев роботизированного склада.
        \item Проведена экспериментальная проверка эффективности, устойчивости и масштабируемости метода.
        \item Сформулированы ограничения метода и направления дальнейшего развития.
    \end{enumerate}
\end{frame}

\begin{frame}
    \frametitle{Ограничения и дальнейшая работа}
    \small
    \begin{columns}[T,onlytextwidth]
        \begin{column}{0.48\textwidth}
            \textbf{Ограничения}
            \begin{itemize}
                \item свойства VCG относятся к отдельному раунду при фиксированных режимах;
                \item платежи и выбор победителей могут быть дороги на крупных масштабах;
                \item качество политики зависит от сценариев обучения;
                \item кластеризация полезна только при достаточном масштабе и локальности графа.
            \end{itemize}
        \end{column}
        \begin{column}{0.48\textwidth}
            \textbf{Дальнейшие направления}
            \begin{itemize}
                \item выбор оркестратора распределения;
                \item приближённые VCG-платежи;
                \item более сложные ограничения и нелинейные стоимости;
                \item перенос политики между конфигурациями;
                \item проверка в более детализированных робототехнических симуляторах.
            \end{itemize}
        \end{column}
    \end{columns}
\end{frame}

\begin{frame}
    \frametitle{Апробация и публикации}
    \small
    Результаты докладывались на:
    \begin{itemize}
        \item 11-й Международной конференции <<Распределенные вычисления и грид-технологии в науке и образовании>> (GRID'2025), Дубна, 2025.
        \item Ежегодной всероссийской научной конференции по проблемам информатики, секция <<Методы машинного обучения и архитектуры интеллектуальных систем>>, Санкт-Петербург, 2026.
    \end{itemize}
    Публикации:
    \begin{itemize}
    	\item Применение механизма аукционов и мультиагентных технологий для распределения ресурсов и задач в распределенной вычислительной системе. в: Physics of Particles and Nuclei. 2026 ; Том 57, № 4. (\textbf{ожидает публикации})
    	\item Алгоритм QMIX для адаптивного управления ресурсами в гетерогенных системах интернета вещей // Вестник Воронежского государственного университета. Серия: Системный анализ и информационные технологии. (\textbf{принято к рецензированию})
    	\item Государственная регистрация программы для ЭВМ: Программа для децентрализованного распределения ресурсов на уровне периферийных вычислений с использованием мультиагентных аукционных механизмов (MA-VCG) – 2026. – RU2026619677 – 06.04.2026 Бюл. № 4.
    \end{itemize}
\end{frame}

\begin{frame}[plain,noframenumbering]
    \begin{center}
        \Huge Спасибо за внимание!\\[4mm]
        \large Готов ответить на вопросы
    \end{center}
\end{frame}
